% Part: first-order-logic
% Chapter: sequent-calculus
% Section: quantifier-rules

\documentclass[../../../include/open-logic-section]{subfiles}

\begin{document}

\olfileid{fol}{seq}{qrl}

\olsection{Regels voor kwantoren}

\subsection{Regels voor $\lforall$}

\begin{defish}
\Axiom$ !A(t), \Gamma \fCenter \Delta$
\RightLabel{\LeftR{\lforall}}
\UnaryInf$ \lforall[x][!A(x)],\Gamma \fCenter \Delta$
\DisplayProof
\hfill
\Axiom$ \Gamma \fCenter \Delta, !A(a) $
\RightLabel{\RightR{\lforall}}
\UnaryInf$ \Gamma \fCenter \Delta, \lforall[x][!A(x)]$
\DisplayProof
\end{defish}

Bij \LeftR{\lforall} is $t$ een gesloten term, dus een term zonder
variabelen. Bij \RightR{\lforall} is $a$~!!a{constant} die nergens in
de onderste sequent van de regel \RightR{\lforall} mag voorkomen. Die
constante moet voor die onderste sequent dus nieuw zijn. We noemen $a$
de \emph{eigenvariabele} van de
\RightR{\forall}-afleidingsstap.\footnote{In deze regel is $a$ eigenlijk
  !!a{constant}. Toch gebruiken we de term
  ``eigenvariabele''. Dat heeft historische redenen.}

\subsection{Regels voor $\lexists$}

\begin{defish}
\Axiom$ !A(a), \Gamma \fCenter \Delta $
\RightLabel{\LeftR{\lexists}}
\UnaryInf$ \lexists[x][!A(x)], \Gamma \fCenter \Delta$
\DisplayProof
\hfill
\Axiom$ \Gamma \fCenter \Delta, !A(t) $
\RightLabel{\RightR{\lexists}}
\UnaryInf$ \Gamma \fCenter \Delta, \lexists[x][!A(x)]$
\DisplayProof
\end{defish}

Ook hier is $t$~een gesloten term. $a$~is !!a{constant} die niet in de
onderste sequent van \LeftR{\lexists} voorkomt. We noemen $a$ de
\emph{eigenvariabele} van de \LeftR{\lexists}-afleidingsstap.

De eigenvariabele moet dus buiten de onderste sequent blijven bij
\RightR{\lforall} en \LeftR{\lexists}. Deze eis heet de
\emph{eigenvariabelevoorwaarde}.

\begin{explain}
We spreken het volgende af. Als $!A$ !!a{formula} is waarin de
!!{variable}~$x$ vrij voorkomt, schrijven we~$!A(x)$. In dezelfde
context is $!A(t)$ kort voor~$\Subst{!A}{t}{x}$. We kunnen de regel
\RightR{\lexists} dus ook zo schrijven:
\begin{prooftree}
  \Axiom$\Gamma \fCenter \Delta, \Subst{!A}{t}{x}$
  \RightLabel{\RightR{\lexists}}
  \UnaryInf$\Gamma \fCenter \Delta, \lexists[x][!A]$
\end{prooftree}
$t$ mag al in~$!A$ voorkomen. Zo kan $!A$ de formule
$\Atom{\Obj P}{t,x}$ zijn. We mogen daarom correct $\Gamma \Sequent
\Delta, \lexists[x][\Atom{\Obj P}{t,x}]$ afleiden uit~$\Gamma \Sequent
\Delta, \Atom{\Obj P}{t,t}$ met~\RightR{\lexists}. Daarbij
``vervangen'' we in de premisse een of
meer voorkomens van~$t$ door de gebonden !!{variable}~$x$. We hoeven
niet alle voorkomens te vervangen. Voor de eigenvariabele in
\RightR{\lforall} en~\LeftR{\lexists} geldt wel een strengere eis: de
!!{constant}~$a$ mag niet in~$!A$ voorkomen. Daarom mogen we niet
$\Gamma \Sequent \Delta, \lforall[x][\Atom{\Obj P}{a,x}]$ afleiden
uit $\Gamma \Sequent \Delta, \Atom{\Obj P}{a,a}$
met~$\RightR{\lforall}$.
\end{explain}

\begin{explain}
Bij \RightR{\lexists} en \LeftR{\lforall} gelden voor de term~$t$ geen
verdere beperkingen. Bij \LeftR{\lexists} en \RightR{\lforall} geldt
wel de eigenvariabelevoorwaarde. Die eist dat de !!{constant}~$a$
nergens buiten~$!A(a)$ in de bovenste sequent voorkomt. De voorwaarde
is nodig om ervoor te zorgen dat het systeem correct is en dus alleen
geldige sequenten !!{derive}s. Zonder deze voorwaarde zouden de
volgende afleidingen zijn toegestaan:
\begin{prooftree}
  \Axiom$!A(a) \fCenter !A(a)$
  \RightLabel{*\LeftR{\lexists}}
  \UnaryInf$\lexists[x][!A(x)] \fCenter !A(a)$
  \RightLabel{\RightR{\lforall}}
  \UnaryInf$\lexists[x][!A(x)] \fCenter \lforall[x][!A(x)]$
  \DisplayProof\bottomAlignProof
  \qquad
  \Axiom$!A(a) \fCenter !A(a)$
  \RightLabel{*\RightR{\lforall}}
  \UnaryInf$!A(a) \fCenter \lforall[x][!A(x)]$
  \RightLabel{\LeftR{\lexists}}
  \UnaryInf$\lexists[x][!A(x)] \fCenter \lforall[x][!A(x)]$
\end{prooftree}
Maar $\lexists[x][!A(x)] \Sequent \lforall[x][!A(x)]$ is niet geldig.
\end{explain}
\end{document}
